\documentclass[10pt,a4paper]{article}
\usepackage[utf8]{inputenc}
\usepackage[T1]{fontenc}
\usepackage[french]{babel}
\usepackage{lmodern}
\usepackage{amsmath,amssymb}
\usepackage[margin=2.2cm]{geometry}
\usepackage{microtype}
\usepackage{booktabs}
\usepackage{hyperref}
\usepackage{pgfplots}
\pgfplotsset{compat=1.18}
\usepackage{listings}
\lstset{basicstyle=\ttfamily\small,frame=single,breaklines=true,
  columns=fullflexible,keepspaces=true}
\hypersetup{colorlinks=true,urlcolor=blue,linkcolor=blue}
\title{Projet MLSF : résolution du jeu Meowdoku\\
  \large Question 6 : conception et évaluation du solveur optimisé}
\author{}
\date{}

\begin{document}
\maketitle
\vspace{-1em}

\section{Choix de conception}

La question 5 a mis en évidence deux résultats : H8, qui privilégie les
littéraux positifs, est adaptée à Meowdoku ; le recalcul d'une heuristique
sur toute la formule peut coûter plus que le gain obtenu dans l'arbre de
recherche. Le mode \texttt{optimized} de \texttt{sat.py} appelle désormais
un moteur distinct, défini dans \texttt{optimized.py}. Le mode
\texttt{dpll} conserve l'algorithme 2 de la question 2. Le nouveau moteur
reste un DPLL complet, sans apprentissage de clauses ni dépendance externe.

\subsection{Propager sans recopier la formule}

Le DPLL initial reconstruit la liste des clauses à chaque affectation.
Dans le moteur optimisé, les clauses sont conservées et les affectations
sont enregistrées dans un tableau de valeurs et une pile chronologique
(\emph{trail}). Une file implicite, repérée par un indice dans cette pile,
contient les affectations dont les conséquences restent à propager.

\paragraph{Clauses binaires.} Une clause $(a\lor b)$ est représentée par
les implications $\neg a\Rightarrow b$ et $\neg b\Rightarrow a$.
Fixer un littéral à vrai ne visite donc que ses conséquences. Cette
représentation convient particulièrement aux nombreuses contraintes
Meowdoku $(\neg x_{i,j}\lor\neg x_{k,l})$ : placer un chat exclut directement
les cases incompatibles.

\paragraph{Clauses d'au moins trois littéraux.} Deux littéraux sont
surveillés. Lorsqu'un littéral surveillé devient faux, on cherche un
remplaçant non faux. S'il n'en existe pas, l'autre surveillance est soit
vraie (clause satisfaite), soit libre (propagation unitaire), soit fausse
(conflit). Les clauses qui ne sont pas touchées ne sont pas parcourues.
Les clauses unitaires initiales sont affectées avant la recherche ; une
clause vide entraîne immédiatement UNSAT.

\paragraph{Retour arrière.} Chaque niveau mémorise sa décision et la
position de début dans le \emph{trail}. Un échec annule seulement les
affectations de ce niveau, puis essaie la polarité opposée. La recherche
est itérative : elle ne dépend pas de la limite de récursion de Python.
Les tautologies et les clauses identiques sont éliminées au prétraitement.
La règle du littéral pur n'est pas recherchée pendant la résolution, pour
éviter un parcours global supplémentaire ; elle n'est pas nécessaire à
la correction de DPLL.

\subsection{Une heuristique adaptée à la forme de la CNF}

Le prétraitement reconnaît une classe syntaxique : toutes les clauses sont
soit entièrement positives, soit des exclusions binaires négatives.
L'encodage Meowdoku de la question 1 appartient à cette classe, mais le
test ne repose ni sur le nom du fichier ni sur une taille de grille.

Pour cette classe, on choisit la clause positive non satisfaite qui a le
moins de candidats encore libres, puis on essaie un candidat à vrai.
Les égalités entre clauses suivent leur ordre d'entrée ; au sein d'une
clause, les candidats sont départagés par le score positif initial
$J^+(x)=\sum_{C\ni x}2^{-|C|}$, puis par le plus petit numéro de variable.
C'est une adaptation de H8 : la priorité aux clauses positives courtes est
actualisée, mais seules les clauses positives sont parcourues. Sur
Meowdoku, il s'agit des clauses imposant au moins un chat par région,
beaucoup moins nombreuses que les exclusions. Ce choix n'est pas
strictement H8d lorsque plusieurs clauses positives se recouvrent.

Pour les autres CNF, le moteur utilise H4d : il compte les occurrences dans
les clauses résiduelles de taille minimale, choisit la variable la plus
fréquente et essaie sa polarité majoritaire. Les égalités sont départagées
par le plus petit numéro de variable, puis par la polarité positive.

\clearpage
\section{Correction, interface et validation}

\subsection{Invariants de correction}

Les implications binaires sont équivalentes à leurs clauses d'origine.
Les surveillances détectent toute clause devenue unitaire ou contradictoire
avant la prochaine décision. Au retour arrière, une variable passe de
fixée à libre : les surveillances n'ont pas à être restaurées, car aucune
nouvelle valeur fausse n'est introduite. La position de propagation est
replacée au début du niveau annulé.

Le branchement explore toujours les deux valeurs d'une variable libre.
Un conflit à la racine, ou l'épuisement des deux branches à chaque niveau,
permet de conclure UNSAT. Pour une CNF générale, SAT n'est renvoyé qu'une
fois toutes les clauses satisfaites. Pour la classe de placement, il
suffit que toutes les clauses positives soient satisfaites : les
implications ont déjà interdit deux valeurs vraies incompatibles, et
compléter les variables restantes à faux préserve les exclusions.
Le modèle retourné affecte toutes les variables déclarées, même absentes
des clauses. Ces arguments ne supposent pas que la grille soit solvable.

\paragraph{Coût et limites algorithmiques.} Le stockage des clauses,
implications, surveillances et affectations est linéaire dans la taille
de la CNF, hors surcoût des objets Python. Le tri et la déduplication sont
effectués une fois. Une propagation ne recopie pas la CNF ; elle visite
les implications et clauses concernées. Une clause longue peut toutefois
être parcourue plusieurs fois. La recherche reste exponentielle dans le
pire cas ; les surveillances réduisent le coût des nœuds, pas nécessairement
leur nombre. H4d parcourt encore la formule à chaque décision.

\subsection{Interface imposée}

\begin{lstlisting}
python sat.py instance.cnf --mode dpll
python sat.py instance.cnf --mode optimized --output-model
\end{lstlisting}
Par défaut, le mode est \texttt{optimized}. La sortie standard contient
\texttt{s SATISFIABLE} ou
\texttt{s UNSATISFIABLE}, et, pour SAT avec \texttt{--output-model}, une
ligne \texttt{v ... 0}. Aucun modèle n'est écrit pour UNSAT.

L'option facultative \texttt{--stats} écrit sur \textbf{stderr} un objet
JSON contenant les durées et les compteurs. Elle ne modifie pas le format
de sortie imposé. Les décisions comptent les nouveaux branchements ; les
retours arrière comptent le passage à la polarité opposée ; les
propagations comptent les nouvelles affectations forcées, unités initiales
comprises. Le compteur de littéraux purs vaut zéro dans le moteur optimisé.
Le temps de l'heuristique exclut le prétraitement et inclut la dernière
vérification de terminaison. L'option \texttt{--seed} fixe les choix
aléatoires du mode DPLL.

\subsection{Vérifications effectuées}

\begin{itemize}
  \item Les tests du projet et les nouveaux tests passent : \textbf{54 tests}.
  \item Le moteur optimisé est comparé à un oracle exhaustif sur
    \textbf{2\,000 CNF aléatoires} de 1 à 9 variables, dont des CNF de
    placement. Chaque modèle SAT est vérifié clause par clause.
  \item Des tests ciblent les clauses vides, formules vides, tautologies,
    doublons, unités contradictoires, déplacements de surveillance suivis
    de retours arrière, ainsi que des recherches et propagations profondes.
  \item Les tests d'interface contrôlent les deux modes, la complétude du
    modèle et la séparation stdout/stderr. Le banc de mesure rejette un
    modèle faux, incomplet, répété ou hors limites et une sortie parasite.
\end{itemize}
Les tests augmentent la confiance dans l'implémentation ; ils ne remplacent
pas les invariants de correction ci-dessus.

\clearpage
\section{Évaluation expérimentale}

\subsection{Protocole reproductible}

\texttt{bench\_q6.py} lance réellement \texttt{sat.py} dans un nouveau
processus par résolution, avec les options \texttt{--output-model} et \texttt{--stats}.
Le délai est fixé à \textbf{60 secondes par processus}, et les lancements
sont séquentiels. Le temps principal inclut le démarrage de Python, les
imports, la lecture, la résolution et l'écriture du résultat. Le temps CPU
interne inclut le prétraitement et la résolution, mais exclut les entrées
et sorties. La vérification indépendante du modèle intervient ensuite,
hors mesure.

Le corpus comprend \textbf{91 instances} : 30 benchmarks publics,
11 tests publics d'interface et de cas limites, 20 grilles Meowdoku SAT
de tailles 8, 10, 12, 14 et 16 (quatre par taille), et 30 grilles à régions
aléatoires (7 SAT, 23 UNSAT). Les 50 grilles et leurs graines sont celles
de la question 5. Les statuts de référence proviennent du fichier public
\texttt{expected.json}, de la construction des grilles SAT et, pour les
régions aléatoires, de l'accord des heuristiques observé à la question 5.

La campagne comporte une exécution par instance et par mode, soit
\textbf{182 résolutions}, avec une graine DPLL dérivée du nom du fichier.
Les mesures sont locales : Python 3.12.13, macOS 15.7.3, architecture ARM64.
Elles ne sont pas des mesures sur la machine Linux x86\_64 de la compétition.
Le CSV \texttt{results/q6\_runs.csv} conserve les résultats détaillés ; le
fichier JSON associé décrit l'environnement et les paramètres.

\subsection{Résultats globaux}

\begin{center}\small
\begin{tabular}{llrrr}
\toprule
Corpus & Mode & Résolus / essais & Temps total (s) & PAR-2 (s) \\
\midrule
Ensemble & dpll & 89/91 & 72.949 & 3.439 \\
Ensemble & optimized & 90/91 & 9.297 & 1.421 \\
Benchmarks publics & dpll & 28/30 & 69.624 & 10.321 \\
Benchmarks publics & optimized & 29/30 & 7.354 & 4.245 \\
Interface et cas limites & dpll & 11/11 & 0.317 & 0.029 \\
Interface et cas limites & optimized & 11/11 & 0.328 & 0.030 \\
Meowdoku SAT & dpll & 20/20 & 1.290 & 0.064 \\
Meowdoku SAT & optimized & 20/20 & 0.687 & 0.034 \\
Régions aléatoires & dpll & 30/30 & 1.718 & 0.057 \\
Régions aléatoires & optimized & 30/30 & 0.928 & 0.031 \\
\bottomrule
\end{tabular}

\end{center}
Le temps total ne porte que sur les résolutions correctes : il doit être
lu avec leur nombre, qui est le premier critère de classement. Le PAR-2
est la moyenne des temps, avec une pénalité de 120 s pour chaque échec
ou dépassement de délai. Il rend ces dépassements visibles dans une
mesure globale ; il n'est pas le critère de classement imposé.

Le mode optimisé obtient 90/91 résolutions correctes, contre 89/91 pour DPLL. La campagne compte 0 erreur(s) de statut, de modèle ou d'interface.
Dépassements de délai (dpll) : \texttt{public-03.cnf}, \texttt{public-06.cnf}.
Dépassements de délai (optimized) : \texttt{public-03.cnf}.
Sur les 89 paires résolues par les deux modes, la moyenne géométrique du rapport DPLL/optimisé vaut 1.81 pour le temps complet et 4.74 pour le temps CPU de résolution.
Le gain n'est pas uniforme : sur \texttt{public-04.cnf}, le mode optimisé est 59.5 fois plus lent en temps complet pour cette graine de DPLL.


\begin{center}
\begin{tikzpicture}
\begin{semilogxaxis}[width=0.90\linewidth,height=5.3cm,
  xlabel={temps complet par instance (s)},ylabel={instances résolues},
  xmin=0.01,xmax=60,ymin=0,ymax=95,grid=major,
  legend pos=south east,legend style={font=\small}]
\addplot[black,thick,mark=none] table[x=wall,y=solved] {generated/cactus-dpll.dat};
\addlegendentry{DPLL initial}
\addplot[blue,thick,mark=none] table[x=wall,y=solved] {generated/cactus-optimized.dat};
\addlegendentry{Optimisé}
\end{semilogxaxis}
\end{tikzpicture}
\end{center}

\clearpage
\section{Interprétation pour Meowdoku}

\begin{center}\footnotesize\setlength{\tabcolsep}{4pt}
\begin{tabular}{rrrrr}
\toprule
Taille & CPU DPLL (ms) & CPU optimisé (ms) & Processus DPLL (ms) & Processus optimisé (ms) \\
\midrule
8 & 2.86 & 0.87 & 32.91 & 29.09 \\
10 & 4.56 & 1.91 & 32.66 & 29.61 \\
12 & 18.43 & 2.69 & 47.66 & 32.12 \\
14 & 47.62 & 4.52 & 77.71 & 34.34 \\
16 & 44.23 & 6.94 & 79.67 & 41.64 \\
\bottomrule
\end{tabular}

\end{center}
Ce tableau donne les temps médians sur les quatre grilles SAT de chaque
taille. Les temps CPU isolent le travail du solveur ; les temps du processus
correspondent davantage à l'expérience de la compétition.

Sur les 50 paires Meowdoku résolues par les deux modes, le gain géométrique est de \textbf{4.2 fois} sur le temps CPU du solveur, contre \textbf{1.5 fois} sur le temps complet. Le démarrage de Python et la lecture de la CNF réduisent donc le gain visible de bout en bout.
Le mode optimisé réussit 50/50 résolutions Meowdoku.


\paragraph{Pourquoi le gain ?} Une décision positive place un chat et
déclenche directement les exclusions binaires de sa ligne, sa colonne,
sa région et son voisinage. La formule n'est plus reconstruite pour chaque
case exclue. Le choix de la région la plus contrainte ne parcourt que les
clauses positives. Le coût d'actualisation reste ainsi limité, alors que
les heuristiques dynamiques de la question 5 recalculaient leurs scores
sur l'ensemble des clauses.

Le gain mesuré combine un nouveau moteur de propagation et une nouvelle
heuristique ; ce protocole ne permet pas d'attribuer une part chiffrée à
chacune. Il compare la solution finale à l'algorithme 2, sans réutiliser
les temps anciens de la question 5 comme s'ils appartenaient à la même
campagne.

\paragraph{Portée des résultats.} Le DPLL de référence n'est exécuté
qu'avec une graine par instance. Une campagne à plusieurs graines serait
nécessaire pour quantifier la variabilité de ses temps. Les mesures très
courtes sont également sensibles à la charge et à la fréquence du
processeur. La campagne actuelle n'est ni une garantie de gain sur toute
CNF, ni une estimation du classement sur les instances privées.

L'absence d'apprentissage de clauses limite les performances sur les
instances exigeant de très nombreux retours arrière. Le mode général H4d
conserve aussi un coût de parcours global par décision. Les grilles à
régions aléatoires du corpus sont majoritairement faciles ; des Meowdoku
UNSAT plus difficiles et des tailles supérieures permettraient d'évaluer
plus largement la robustesse du solveur.

\section{Reproduire les résultats}

Depuis la racine du projet, avec Python 3.12 :
\begin{lstlisting}
python -m unittest discover -v
python bench_q6.py run --timeout 60
python bench_q6.py report
\end{lstlisting}
La première commande vérifie la correction. La deuxième recrée les
instances Meowdoku et le CSV des mesures ; la troisième affiche la
synthèse et régénère les tableaux et les données des courbes. Pour
répéter uniquement les 50 grilles Meowdoku sans écraser la campagne
principale :
\begin{lstlisting}
python bench_q6.py run --suite meowdoku --repeats 3 \
  --csv results/q6_meow_repeats.csv
python bench_q6.py report --csv results/q6_meow_repeats.csv
\end{lstlisting}
Les répétitions conservent la même graine : elles mesurent la variabilité
temporelle, pas celle des choix aléatoires. Les durées et les compteurs
d'une seule résolution sont accessibles avec :
\begin{lstlisting}
python sat.py instance.cnf --mode optimized --output-model --stats
\end{lstlisting}

\end{document}
